\chapter*{Conclusiones}

De nuevo unos sabios consejos tomados de \cite{guidedissertation}: \\
Here you will bring together the work of the dissertation by showing how the
initial research plan has been addressed in such a way that conclusions may be
formed from the evidence of the dissertation. No new material or references
should be placed here. The conclusions should make a statement on the extent
to which each of the aims and objectives has been met. You should bring back
your research questions and state clearly your understanding of those questions.
Be careful not to make claims that are not substantiated from the evidence you
have presented in earlier chapters.

If you are undertaking a company project based around a business issue do not
confuse recommendations for the company with conclusions. If you want to
include a list of recommendations then do so in a separate short chapter.
\emph{The conclusions address the wider understanding of the issue you have
been studying.}

You should include a short sub section on any suggestions for further research
for colleagues who might wish to undertake research in this area in the future.
There should also be a short statement of the limitations of the research. Often
as a single case study or limited range of companies you can not really claim
that your research holds for all companies. However, by adopting a rigorous
approach to your literature review and methods which have validity and can be
repeated you can make a reasonable but limited claim that your conclusions
should be taken seriously.

